% ============ 12 高阶直像; Leray 谱序列 (原书 pp.81-83) ============
\chapter{高阶直像；Leray 谱序列}\label{sec:12}
\index{高阶直像}
\index{Leray 谱序列}

\subsection*{高阶直像}

设 $\varphi \colon Y \to X$ 是簇(或概形)的态射.回忆对 $Y_{\mathrm{et}}$ 上的层 $F$，我们把 $\varphi_{*}F$ 定义为 $X_{\mathrm{et}}$ 上满足
\[
\Gamma(U, \varphi_{*}F) = \Gamma(U_Y, F), \qquad U_Y \overset{\mathrm{def}}{=} U \times_{X} Y
\]
的层.

函子 $\varphi_{*} \colon \mathrm{Sh}(Y_{\mathrm{et}}) \to \mathrm{Sh}(X_{\mathrm{et}})$ 左正合，因此可以考虑它的右导出函子 $R^{r}\varphi_{*}$.我们称层 $R^{r}\varphi_{*}F$ 为 $F$ 的高阶直像.

\begin{proposition}\label{prop:12.1}
对任何 $\varphi \colon Y \to X$ 与 $Y_{\mathrm{et}}$ 上的层 $F$，$R^{r}\varphi_{*}F$ 是与预层 $U \mapsto H^{r}(U_Y, F)$ 相伴的 $X_{\mathrm{et}}$ 上的层.
\end{proposition}

\begin{proof}
设 $\varphi_{p}$ 是 $Y_{\mathrm{et}}$ 上的预层 $P$ 送到 $X_{\mathrm{et}}$ 上的预层 $U \mapsto \Gamma(U_Y, P)$ 的函子 $\mathrm{PreSh}(Y_{\mathrm{et}}) \to \mathrm{PreSh}(X_{\mathrm{et}})$---它显然正合.由 $\varphi_{*}$ 的定义，图
\[
\begin{tikzcd}
\mathrm{PreSh}(Y_{\mathrm{et}}) \arrow[r, "\varphi_{p}"] & \mathrm{PreSh}(X_{\mathrm{et}}) \arrow[d, "a"] \\
\mathrm{Sh}(Y_{\mathrm{et}}) \arrow[u, "i"] \arrow[r, "\varphi_{*}"'] & \mathrm{Sh}(X_{\mathrm{et}})
\end{tikzcd}
\]
交换($i$ 是``把层看作预层''的函子).设 $F \to \mathcal{I}^{\bullet}$ 是 $F$ 的内射分解.则因为 $a$ 与 $\varphi_{p}$ 正合，
\[
R^{r}\varphi_{*}F \overset{\mathrm{def}}{=} H^{r}(\varphi_{*}\mathcal{I}) = H^{r}(a \circ \varphi_{p} \circ i \, \mathcal{I}^{\bullet}) = a \circ \varphi_{p}(H^{r}(i\mathcal{I}^{\bullet})).
\]
如我们已经注意到的(10.4 的证明)，$H^{r}(i\mathcal{I}^{\bullet})$ 是预层 $U \mapsto H^{r}(U, F)$，从而 $\varphi_{p}(H^{r}(i\mathcal{I}^{\bullet}))$ 是预层 $U \mapsto H^{r}(U_Y, F)$.
\end{proof}

\begin{corollary}\label{cor:12.2}
$R^{r}\varphi_{*}F$ 在 $\bar{x} \to X$ 处的茎是 $\varinjlim H^{r}(U_Y, F)$，其中极限遍历 $\bar{x}$ 的所有平展邻域 $(U, u)$.
\end{corollary}

\begin{proof}
按定义，$\varinjlim H^{r}(U_Y, F)$ 是 $\varphi_{p}(R^{r}i(F))$ 的茎，由定理~\ref{thm:7.15}，它等于 $a\varphi_{p}(R^{r}i(F))$ 的茎.
\end{proof}

\begin{example}\label{ex:12.3}
若 $\varphi \colon Y \hookrightarrow X$ 是闭浸入，则 $\varphi_{*}$ 正合，从而对 $r > 0$，$R^{r}\varphi_{*}F = 0$.更一般地，若 $\varphi \colon Y \to X$ 是有限映射，则对 $r > 0$，$R^{r}\varphi_{*}F = 0$(同样因为 $\varphi_{*}$ 正合，8.4).
\end{example}

\begin{example}\label{ex:12.4}
假定 $X$ 连通且正规，设 $g \colon \xi \hookrightarrow X$ 是 $X$ 的一般点的包含.则
\[
(R^{r}g_{*}F)_{\bar{x}} = H^{r}(\operatorname{Spec} K_{\bar{x}}, F),
\]
其中 $K_{\bar{x}}$ 是 $\mathcal{O}_{X, \bar{x}}$ 的分式域.此外，在这种情况下，$g_{*}$ 把 $\xi$ 上的常值层送到常值层.我们把它显式化.

设 $K^{\mathrm{sep}}$ 是 $K$ 的可分闭包，$G = \operatorname{Gal}(K^{\mathrm{sep}}/K)$.设 $M = M_F$，即与 $F$ 对应的 $G$-模(如 \S 6 中).则 $F$ 常值，当且仅当 $G$ 平凡作用于 $M$；$F$ 局部常值，当且仅当 $G$ 在 $M$ 上的作用穿过有限商.

映射 $\operatorname{Spec} K^{\mathrm{sep}} \to \operatorname{Spec} K \to X$ 是 $X$ 的几何点，记为 $\bar{\xi}$.严格局部环 $\mathcal{O}_{X, \bar{\xi}}$ 是 $K^{\mathrm{sep}}$，因为 $X$ 在 $K^{\mathrm{sep}}$ 中 $K$ 的任何有限扩张 $L$ 中的正规化会在 $X$ 的某个非空开子集上平展.于是，若 $r = 0$，$(R^{r}g_{*}F)_{\bar{\xi}} = M$，否则为 0.

一般地，$K_{\bar{x}}$ 是 $K^{\mathrm{sep}}$ 中所有这样的有限扩张 $L$ 的并：$X$ 在 $L$ 中的正规化在位于 $x$ 之上的某个点处非分歧.于是 $(R^{r}g_{*}F)_{\bar{x}} = H^{r}(H, M)$，其中 $H = \operatorname{Gal}(K^{\mathrm{sep}}/K_{\bar{x}})$.

例如，设 $X = \operatorname{Spec} A$，$A$ 是 Dedekind 整环.设 $\widetilde{A}$ 是 $A$ 在 $K^{\mathrm{sep}}$ 中的整闭包.$X$ 的闭点 $x$ 是 $A$ 的非零素理想 $\mathfrak{p}$，而位于 $\mathfrak{p}$ 之上的 $\widetilde{A}$ 的素理想 $\widetilde{\mathfrak{p}}$ 的选取确定 $X$ 的几何点 $\bar{x} \to x \to X$.在这种情况下，$K_{\bar{x}} = (K^{\mathrm{sep}})^{I(\widetilde{\mathfrak{p}})}$，其中 $I(\widetilde{\mathfrak{p}}) \subset G$ 是 $\widetilde{\mathfrak{p}}$ 的惯性群.于是 $(R^{r}g_{*}F)_{\bar{x}} = H^{r}(I(\widetilde{\mathfrak{p}}), M)$.
\end{example}

\begin{example}\label{ex:12.5}
假定 $X$ 整，但不一定正规.则 $g \colon \xi \to X$ 分解为 $g \colon \xi \to \widetilde{X} \to X$，其中 $\widetilde{X}$ 是 $X$ 在 $\xi$ 中的正规化.例如：
\[
X = \operatorname{Spec} A; \quad A \text{ 整环};
\qquad
\xi = \operatorname{Spec} K; \quad K \text{ 是 } A \text{ 的分式域};
\qquad
\widetilde{X} = \operatorname{Spec} \widetilde{A}; \quad \widetilde{A} \text{ 是 } A \text{ 在 } K \text{ 中的整闭包}.
\]
由于 $\widetilde{X} \to X$ 有限，这表明 $g$ 分解为上面两个例子中考虑过的类型映射的复合.

对映射 $\widetilde{X} \to X$，常值层的直像不一定是常值的.考虑例如映射
\[
\varphi \colon \mathbb{A}^{1} \to \{Y^{2} = X^{3} + X^{2}\}, \qquad t \mapsto (t^{2} - 1, t(t^{2} - 1)),
\]
从 $\mathbb{A}^{1}$ 到结点三次曲线.映射是一一对应的，除了两个点 $\pm 1$ 映到曲线上的点 $(0, 0)$.因此，对 $\mathbb{A}^{1}$ 上的常值层 $\Lambda$，除非 $x = (0, 0)$(此时它是 $\Lambda \oplus \Lambda$), $(\varphi_{*}\Lambda)_{\bar{x}} = \Lambda$.
\end{example}

\begin{example}\label{ex:12.6}
设 $X$ 是整概形，考虑 $X$ 的点 $z$ 的包含 $i \colon z \to X$($z$ 不一定是一般点或闭点).设 $Z$ 是 $z$ 在 $X$ 中的闭包，则 $Z$ 不可约且 $z$ 是它的一般点.则 $i$ 分解为
\[
z \to \widetilde{Z} \to Z \hookrightarrow X,
\]
它是最后三个例子中考虑过的类型的映射的复合.
\end{example}

\subsection*{Leray 谱序列}

\begin{theorem}\label{thm:12.7}[Leray 谱序列]
设 $\varphi \colon Y \to X$ 是簇(或概形)的态射.对 $Y_{\mathrm{et}}$ 上的任意层 $F$，存在谱序列
\[
H^{r}(X_{\mathrm{et}}, R^{s}\varphi_{*}F) \Rightarrow H^{r+s}(Y_{\mathrm{et}}, F).
\]
\end{theorem}

\begin{proof}
函子 $\varphi_{*} \colon \mathrm{Sh}(Y_{\mathrm{et}}) \to \mathrm{Sh}(X_{\mathrm{et}})$ 与 $\Gamma(X, \cdot) \colon \mathrm{Sh}(X_{\mathrm{et}})$ 都左正合，且它们的复合是 $\Gamma(Y, \cdot)$.由于 $\varphi_{*}$ 保持内射对象，这个定理是定理~\ref{thm:10.7} 的特殊情形.
\end{proof}

原始的 Leray 谱序列与上述序列有完全相同的形式：对拓扑空间的任何连续映射 $\varphi \colon Y \to X$ 与 $Y$ 上的层 $F$，存在谱序列
\[
H^{r}(X, R^{s}\varphi_{*}F) \Rightarrow H^{r+s}(Y, F).
\]

\begin{example}\label{ex:12.8}
设 $X$ 是域 $k$ 上的簇，$\varphi \colon X \to P$ 是把 $X$ 映到点 $P = \operatorname{Specm}(k)$ 的映射.设 $\bar{k}$ 是 $k$ 的可分闭包，$\bar{X}$ 是由 $X$ 基变换得到的 $\bar{k}$ 上的簇.设 $\Gamma = \operatorname{Gal}(\bar{k}/k)$.当把 $\mathrm{Sh}(P_{\mathrm{et}})$ 等同于离散 $\Gamma$-模范畴 $\mathrm{Mod}_{\Gamma}$ 时，$\varphi_{*}$ 等同于函子
\[
\mathrm{Sh}(X_{\mathrm{et}}) \to \mathrm{Mod}_{\Gamma}, \qquad F \mapsto F(\bar{X}) \overset{\mathrm{def}}{=} \varinjlim_{k'} F(X_{k'})
\]
(极限遍历在 $k$ 上有限的 $\bar{k}$ 的子域 $k'$).于是，在这种情况下，Leray 谱序列变为
\[
H^{r}(\Gamma, H^{s}(\bar{X}_{\mathrm{et}}, F)) \Rightarrow H^{r+s}(X, F),
\]
其中
\[
H^{s}(\bar{X}, F) \overset{\mathrm{def}}{=} \varinjlim_{k'} H^{s}(X_{k'}, F|_{X_{k'}}) = H^{s}(\bar{X}, \varphi'^{*}F), \qquad \varphi' \colon \bar{X} \to X.
\]
\end{example}
